% =====================================================================
\chapter{Introdução}\label{cap:introducao}
% =====================================================================

O ambiente corporativo brasileiro é composto por empresas com diferentes
estruturas de propriedade e de governança. Entre elas, destacam-se as empresas
familiares, que representam uma parcela relevante das organizações e possuem
características particulares relacionadas à preservação do controle, sucessão e
continuidade do negócio ao longo das gerações \cite{pestana2021}. Em contraste,
empresas com estruturas acionárias mais dispersas tendem a apresentar modelos de
gestão mais profissionalizados e menor influência de um grupo controlador
específico. Neste cenário, as decisões de estrutura de capital, isto é, a
combinação entre capital próprio e de terceiros, representam um dilema central da
gestão financeira. Teorias clássicas como o \textit{Trade-Off}, que explica a
escolha ótima de endividamento com base no equilíbrio entre benefícios fiscais e
custos de falência, e a \textit{Pecking Order}, que propõe uma hierarquia de
financiamento baseada na preferência por recursos internos, oferecem
interpretações racionais para essas decisões, assumindo como objetivo primordial a
maximização do valor da firma \cite{myers1984,myersmajluf1984}.

Contudo, a aplicação direta destes modelos torna-se complexa no contexto de
empresas familiares. A literatura sobre governança em firmas familiares introduz o
conceito de riqueza socioemocional (\textit{socioemotional wealth} -- SEW), que
representa a utilidade não financeira que a família extrai do negócio
\cite{gomezmejia2007}. Essa
``riqueza'' abrange a manutenção do controle, a perpetuidade do legado e a
identidade familiar atrelada à organização, objetivos que podem se sobrepor à pura
maximização de resultados econômicos \cite{santos2018}.

Essa busca pela preservação do SEW colide diretamente com os riscos inerentes ao
endividamento. Contratos de dívida, especialmente os de longo prazo, frequentemente
impõem cláusulas restritivas (\textit{covenants}) que servem como mecanismo de
disciplina, mas que também representam uma potencial ameaça à autonomia gerencial
\cite{costa2018}. A violação dessas cláusulas pode resultar em penalidades como a
renegociação compulsória das obrigações financeiras ou até mesmo o vencimento
antecipado da dívida. Nesse contexto, verificou-se que empresas com níveis mais
elevados de endividamento podem apresentar incentivos ao gerenciamento de
resultados, buscando reduzir a probabilidade de descumprimento dos
\textit{covenants} e os custos associados à sua violação \cite{costa2018}. Para
uma família controladora, a mera possibilidade de intervenção de credores
representa uma ameaça direta ao seu poder decisório.

Assim, a aversão à interferência de credores, que ameaça o SEW, pode levar a uma
política de endividamento mais conservadora. Como o
setor de atuação também influencia a estrutura de capital \cite{oliveira2020}, a
comparação entre empresas familiares e não familiares precisa reduzir o risco de
que diferenças setoriais sejam confundidas com o efeito da estrutura de controle. Por
isso, esta pesquisa concentra a análise nas companhias de um único setor, o de
saúde da B3 S.A. -- Brasil, Bolsa, Balcão (B3).
\verificar{conferir em Oliveira e Kayo (2020) a passagem sobre o efeito do setor}

Diante do exposto, emerge o seguinte problema de pesquisa: a estrutura de controle
familiar está associada a diferentes níveis de endividamento quando comparada à
estrutura de controle não familiar nas companhias abertas brasileiras do setor de
saúde? A resposta a esta questão pode subsidiar investidores, credores e os
próprios gestores familiares na formulação de estratégias financeiras alinhadas a
seus objetivos econômicos e não econômicos.

O objetivo geral deste trabalho é comparar o nível e a composição do endividamento
de empresas familiares e não familiares do setor de saúde da B3 no período de 2018
a 2025. Como objetivos específicos, busca-se: (i) classificar as companhias quanto
ao controle familiar a partir de critérios da literatura; (ii) comparar os grupos
quanto ao endividamento oneroso, ao endividamento total, à capacidade de pagamento
e, de forma complementar e exploratória, ao prazo da dívida; e (iii) verificar se as diferenças se mantêm em subperíodos e
sob critérios alternativos de classificação.

Nesse contexto, justifica-se a realização deste estudo pela relevância prática e
teórica do tema. Apesar da extensa literatura sobre estrutura de capital e
empresas familiares, representada por estudos como \citeonline{kayo2018},
\citeonline{wildner2020} e, em empresas portuguesas, \citeonline{pestana2021}, os
resultados acerca da
influência da estrutura de controle familiar sobre o endividamento corporativo
permanecem inconclusivos. Enquanto alguns estudos apontam para uma postura mais
conservadora e menores níveis de alavancagem em empresas familiares, outros
identificam comportamentos distintos em função de características específicas de
governança, do ambiente institucional e das particularidades de cada mercado. Além
disso, esta pesquisa examina um período recente, que inclui o choque da pandemia de
COVID-19, concentra a comparação em um único setor e verifica se os resultados se
mantêm sob diferentes critérios de classificação de empresa familiar, aspecto em
que os estudos nacionais divergem. Com isso, busca oferecer uma contribuição
empírica adicional para a literatura de finanças corporativas e de governança em
empresas familiares. Do ponto de vista gerencial, os resultados podem auxiliar
gestores, credores e investidores a compreender de forma mais clara como a
presença da família no controle afeta o apetite ao endividamento e,
consequentemente, o risco financeiro da organização.

Além desta introdução, o trabalho está organizado da seguinte forma: o
\autoref{cap:referencial} apresenta o referencial teórico e a hipótese de
pesquisa; o \autoref{cap:metodologia} descreve a amostra, a classificação das
empresas, os indicadores e os procedimentos estatísticos; o
\autoref{cap:resultados} apresenta e discute os resultados; e o
\autoref{cap:consideracoes} traz as considerações finais.
